\chapter{Os dados: séries temporais de poços e o conjunto 3W}
\label{cap:dados}

\begin{objetivos}
Ao final deste capítulo, você será capaz de:
\begin{itemize}[leftmargin=*,itemsep=0.2em]
  \item formalizar uma série temporal multivariada de poço e definir janela, passo e
        rótulo;
  \item descrever a estrutura do conjunto 3W nas versões 1.0 e 2.0, incluindo a
        distinção entre instâncias reais, simuladas e desenhadas à mão;
  \item explicar a rotulação em três estados --- normal, transiente e anômalo --- e
        distinguir rótulo por amostra de rótulo por instância;
  \item identificar as principais armadilhas do conjunto: desbalanceamento, valores
        ausentes, sensores congelados, vazamento temporal e correlação espúria entre
        anomalia e perda de telemetria;
  \item projetar um protocolo de particionamento que evite vazamento entre treino e
        teste.
\end{itemize}
\end{objetivos}

\section{Formalização}

\begin{definicao}{Série temporal multivariada de poço}
Uma série temporal multivariada de poço é uma função
$\vect{x} : \mathcal{T} \to \R^{S}$, onde $\mathcal{T} \subset \mathbb{Z}$ é um
conjunto de instantes amostrados a intervalo regular $\Delta t$ e $S$ é o número de
sensores. Escrevemos $\vect{x}_t = (x_t^{(1)}, \dots, x_t^{(S)})^{\top}$ para a leitura
conjunta no instante $t$, e $x_t^{(s)}$ para a leitura do sensor $s$.
\end{definicao}

No conjunto 3W, $\Delta t = \SI{1}{\second}$ e uma instância típica cobre de algumas
horas a alguns dias. O subconjunto de sensores usado ao longo deste livro é, salvo
menção em contrário, $S = 4$, com os canais P-PDG, T-PDG, P-TPT e T-TPT descritos na
\cref{sec:dominio-sensores}.

\begin{definicao}{Janela deslizante}
Dada uma série $\vect{x}$, um comprimento de janela $T$ e um passo $h$, a
\termo{janela deslizante}{sliding window} de índice $i$ é a matriz
\[
  \mat{X}_i = \big[\, \vect{x}_{t_0 + ih},\; \vect{x}_{t_0 + ih + 1},\;
  \dots,\; \vect{x}_{t_0 + ih + T - 1} \,\big]^{\top} \in \R^{T \times S}.
\]
Quando $h < T$ as janelas se sobrepõem; quando $h = T$ elas são disjuntas.
\end{definicao}

A janela é a unidade de entrada de praticamente todos os modelos deste livro, e sua
parametrização varia de forma significativa entre eles. A \cref{tab:dados-janelas}
consolida as escolhas, o que é útil como referência ao longo da leitura.

\begin{table}[htbp]
\centering
\caption{Parametrização de janela adotada em cada capítulo deste livro.}
\label{tab:dados-janelas}
\small
\begin{tabular}{llll}
\toprule
\textbf{Capítulo} & \textbf{$T$ (amostras)} & \textbf{$S$} & \textbf{Passo $h$} \\
\midrule
\cref{cap:dual}         & 500                & 5 & contínuo (tempo real) \\
\cref{cap:binarios,cap:closedworld} & 150 a 7500 (varrido) & 8 & sobreposto \\
\cref{cap:mundoaberto}  & 200                & 4 & sobreposto \\
\cref{cap:segmentacao,cap:mahalanobis} & 200 & 4 & 100 (\SI{5}{\second} de avanço) \\
\cref{cap:agente}       & segmento variável  & 4 & --- (segmentos rotulados) \\
\bottomrule
\end{tabular}
\end{table}

\begin{destaque}[A escolha de $T$ não é livre]
O comprimento da janela impõe um compromisso duplo. Janelas curtas permitem resposta
rápida, mas não contêm informação suficiente para eventos lentos --- uma janela de
\SI{150}{\second} não pode, em princípio, distinguir incrustação (classe 7) de
restrição rápida (classe 6). Janelas longas capturam eventos lentos, mas atrasam a
detecção e diluem transientes curtos na média. O \cref{cap:closedworld} quantifica
esse compromisso com uma varredura de catorze valores de $T$.
\end{destaque}

\section{O conjunto 3W}

\subsection{Origem e propósito}

O 3W --- o nome registra que o conjunto reúne instâncias de \emph{três} procedências
distintas (reais, simuladas e desenhadas à mão) de eventos indesejáveis que ocorrem em
poços (\ing{wells}) de petróleo \citep{Petrobras3W} --- foi publicado pela Petrobras
com o objetivo explícito de dar à comunidade acadêmica um banco de dados realista para
detecção e classificação de eventos \citep{VARGAS2019}. Antes dele, a literatura da
área trabalhava quase exclusivamente com dados simulados ou proprietários, o que
tornava impossível comparar métodos.

\subsection{Versão 1.0}

A versão original contém \num{1984} instâncias distribuídas entre a classe normal e
oito classes de anomalia. Cada instância é um arquivo CSV com uma coluna de tempo,
colunas de sensores e uma coluna de rótulo por amostra.

\subsection{Versão 2.0}

A versão 2.0 \citep{VazVargas2026_3w2} amplia o conjunto para \num{2228} instâncias e
\num{10} classes (normal mais nove anomalias), revisa a rotulação de instâncias da
versão anterior e acrescenta metadados de poço. Os trabalhos discutidos nos
\cref{cap:dual,cap:binarios,cap:mundoaberto} foram desenvolvidos sobre a versão 1.0; os
dos \cref{cap:segmentacao,cap:mahalanobis,cap:agente} incorporam a versão 2.0.

\begin{destaque}[Comparabilidade entre versões]
Números de acurácia obtidos em versões diferentes do 3W \emph{não são diretamente
comparáveis}. A revisão de rótulos da versão 2.0 alterou a fronteira entre normal e
transiente em várias instâncias, e a inclusão de novas instâncias mudou a distribuição
de classes. Sempre que este livro compara resultados de capítulos distintos, a versão
utilizada é declarada.
\end{destaque}

\subsection{Três procedências}

Cada instância do 3W pertence a uma de três procedências, identificada pelo prefixo do
nome do arquivo:

\begin{description}[leftmargin=0pt,style=nextline,itemsep=0.5em]

\item[Real (\texttt{WELL-\dots}).] Registro efetivo de um poço em operação, com todos
os defeitos que isso implica: ruído, lacunas, deriva, comportamento idiossincrático.
É a única procedência que sustenta conclusões sobre desempenho em campo.

\item[Simulado (\texttt{SIMULATED-\dots}).] Gerado por simulador de escoamento
multifásico. Tem a vantagem de cobrir condições raras e a desvantagem de refletir as
hipóteses do simulador. Útil para aumentar a amostra de classes escassas.

\item[Desenhado à mão (\texttt{DRAWN-\dots}).] Construído manualmente por especialistas
para ilustrar a assinatura canônica de um evento. \emph{Não deve ser usado para
avaliação}: são exemplos idealizados, e um modelo avaliado sobre eles reporta um
desempenho que não se sustenta em campo.

\end{description}

\begin{destaque}[Uma decisão de projeto recorrente]
O trabalho descrito no \cref{cap:agente} adota a seguinte política, que recomendamos:
usar exclusivamente instâncias reais; recorrer a simuladas apenas quando a classe tem
menos de trinta instâncias reais; e excluir integralmente as desenhadas à mão. Essa
política reduz o tamanho do conjunto e piora os números publicados --- e é exatamente
por isso que ela é correta.
\end{destaque}

\section{A rotulação}

\subsection{Três estados por amostra}

O 3W rotula \emph{cada amostra} de cada instância com um de três estados:

\begin{center}
\small
\begin{tabular}{cll}
\toprule
\textbf{Código} & \textbf{Estado} & \textbf{Significado} \\
\midrule
$0$        & Normal    & Operação estável, anterior ao evento. \\
$100 + k$  & Transiente& Transição em curso para a anomalia de classe $k$. \\
$k$        & Anômalo   & Regime anômalo estabelecido da classe $k$. \\
\bottomrule
\end{tabular}
\end{center}

O estado transiente é a contribuição mais original --- e mais problemática --- do
esquema de rotulação. Ele reconhece que a passagem de normal para anômalo não é
instantânea, e que existe um intervalo em que o evento já começou mas ainda não se
manifestou plenamente. Do ponto de vista operacional, \emph{o transiente é exatamente
onde a detecção tem valor}: detectar depois que o regime anômalo se estabeleceu é tarde
demais.

\begin{destaque}[O que fazer com o transiente]
Três políticas convivem na literatura e neste livro:
\begin{enumerate}[leftmargin=*,itemsep=0.2em]
  \item \textbf{Descartar} as amostras transientes do treinamento. Produz classes mais
        puras e números melhores, ao custo de tornar o modelo cego justamente à fase
        que interessa.
  \item \textbf{Tratar como anômalo}. Favorece detecção precoce e aumenta o falso
        positivo.
  \item \textbf{Tratar como classe própria}. Triplica o número de classes e agrava o
        desbalanceamento.
\end{enumerate}
Os trabalhos deste livro adotam predominantemente a segunda política. A métrica SoftED
do \cref{cap:segmentacao} foi introduzida precisamente para avaliar a detecção com
tolerância temporal, tornando a questão menos binária.
\end{destaque}

\subsection{Rótulo por instância}

Além do rótulo por amostra, cada instância possui um rótulo global --- a classe do
evento que ela contém. Essa distinção gera duas famílias de métricas que \emph{não
podem ser confundidas}:

\begin{itemize}[leftmargin=*,itemsep=0.3em]
  \item \textbf{Métrica por amostra}: fração de amostras corretamente rotuladas.
        Dominada pelas instâncias mais longas e pela classe normal.
  \item \textbf{Métrica por instância}: fração de instâncias em que o evento foi
        corretamente identificado. É a que corresponde à pergunta operacional.
\end{itemize}

\begin{exemplo}{A diferença em números}
No \cref{cap:mundoaberto}, o classificador binário da classe 6 apresenta acurácia média
por amostra de $0{,}75 \pm 0{,}13$ e acurácia média por instância de
$0{,}34 \pm 0{,}42$. O desvio-padrão enorme na segunda métrica revela o que a primeira
esconde: o modelo acerta quase todas as amostras de algumas instâncias e erra quase
todas as de outras. Uma média de $0{,}75$ por amostra é compatível com ``funciona bem
em metade dos poços e não funciona na outra metade''.
\end{exemplo}

\section{As armadilhas}

\subsection{O estado das válvulas muda a estrutura de correlação}

Antes das armadilhas de rotulação, uma armadilha de estacionariedade. As variáveis
observadas não guardam entre si uma relação fixa: a estrutura de correlação depende da
configuração das válvulas no momento da medição.

\begin{figure}[htbp]
\centering
\includegraphics[width=0.62\textwidth]{correlation.png}
\caption{Matrizes de correlação entre os sinais de pressão e temperatura (PT, TPT e PDG),
calculadas separadamente para dados do 3W e para dados de operação com a válvula M2
aberta e fechada. Vermelho indica correlação positiva; azul, negativa. Adaptado de
\citet{aranha2023system}.}
\label{fig:dados-correlacao}
\end{figure}

A \cref{fig:dados-correlacao} mostra o efeito. Pares de sensores que se correlacionam
positivamente com a M2 aberta podem correlacionar-se de forma fraca ou invertida com a
M2 fechada, porque o caminho do fluido mudou. Um modelo treinado sobre a mistura dos dois
regimes aprende uma correlação média que não descreve nenhum dos dois.

\begin{destaque}[A consequência prática]
É por isso que o sistema do \cref{cap:dual} suspende a análise durante manobras de
válvula e retoma apenas após a estabilização, e é por isso que ele mantém modelos
distintos por configuração. Não é excesso de zelo: sem essa segmentação, toda manobra
programada vira um alarme.

O mesmo cuidado não é possível no 3W, que não distribui o estado das válvulas. Essa
ausência é uma limitação do conjunto, e não do método.
\end{destaque}

\subsection{Desbalanceamento}

A classe normal domina o conjunto por ordens de magnitude ao nível de amostra, e as
classes de anomalia são desbalanceadas entre si --- a classe 9 tem uma pequena fração
das instâncias da classe 7. As consequências são conhecidas e as contramedidas também:
métricas balanceadas (\cref{cap:aprendizado}), amostragem estratificada, ponderação de
classe na função de perda e, no caso do ensemble binário do \cref{cap:binarios},
subamostragem da classe negativa.

\subsection{Valores ausentes e sensores congelados}

Instâncias reais contêm lacunas. As estratégias usadas neste livro são:

\begin{itemize}[leftmargin=*,itemsep=0.3em]
  \item preenchimento com o último valor válido (\ing{forward fill}), adequado quando a
        lacuna é curta;
  \item preenchimento com o próximo valor válido (\ing{backward fill}), para lacunas no
        início da série;
  \item preenchimento pela média da janela, usado no \cref{cap:segmentacao};
  \item descarte da instância, quando a fração ausente excede um limiar.
\end{itemize}

Um sensor congelado não gera valor ausente --- gera um valor constante, que passa por
todos os filtros de qualidade baseados em nulos. Detectá-lo exige verificar
explicitamente $\std(x^{(s)}_{t:t+T}) < \epsilon$.

\subsection{A correlação espúria mais perigosa do 3W}
\label{sec:dados-telemetria}

Este é o achado mais desconfortável reportado neste livro, e ele aparece no
\cref{cap:resultados-agente}. Em várias classes do 3W --- notadamente 2, 6, 7 e 8 ---
o instante de início da anomalia coincide, com frequência muito acima do acaso, com a
perda dos medidores de fundo: P-PDG e T-PDG tornam-se ausentes ou zerados.

A explicação provável é banal e não tem nada a ver com física de poços. Muitos eventos
indesejáveis são acompanhados de intervenção, e intervenção frequentemente implica
reconfiguração da aquisição. O rótulo de anomalia e a interrupção da telemetria têm uma
causa comum.

\begin{destaque}[Por que isso importa]
Um modelo treinado nesses dados pode aprender ``se os sensores de fundo sumiram, é
anomalia''. Essa regra tem alto desempenho no conjunto de teste e valor operacional
próximo de zero. Pior: ela é invisível às métricas usuais.

Foi um modelo de linguagem, ao ser solicitado a \emph{descrever em texto} o que
observava, que tornou o padrão explícito: em quatro das sete classes analisadas, a
descrição gerada continha a expressão ``perda de telemetria de fundo''. O episódio é o
argumento mais forte deste livro a favor da explicabilidade --- não porque o modelo de
linguagem classifique melhor, mas porque ele \emph{fala}, e ao falar expõe o que o
classificador silencioso teria escondido.
\end{destaque}

\subsection{Vazamento temporal e vazamento entre poços}

Duas formas de vazamento comprometem avaliações neste domínio.

\begin{description}[leftmargin=0pt,style=nextline,itemsep=0.4em]

\item[Vazamento temporal.] Ocorre quando janelas sobrepostas de uma mesma instância
são distribuídas entre treino e teste. Como janelas adjacentes compartilham a maior
parte de suas amostras, o modelo essencialmente memoriza. A contramedida é a
\emph{divisão temporal}: os primeiros 80\,\% de cada instância para treino, os últimos
20\,\% para teste, sem sobreposição na fronteira. É o protocolo adotado no
\cref{cap:closedworld}.

\item[Vazamento entre poços.] Ocorre quando janelas de um mesmo poço aparecem em treino
e teste. O modelo aprende as idiossincrasias daquele poço --- nível de pressão
característico, ruído do sensor --- em vez da assinatura do evento. A contramedida é a
\emph{divisão por poço}, mais rigorosa e mais realista. O \cref{cap:segmentacao} reporta
resultados por poço individual justamente para expor a variabilidade que uma média
global esconde.

\end{description}

\begin{exemplo}{Quanto custa fazer certo}
No \cref{cap:segmentacao}, o desempenho da U-Net varia de $0{,}720$ de acurácia no poço
F1 a $0{,}975$ no poço C2. A média global, $0{,}863 \pm 0{,}082$, só é informativa
porque vem acompanhada do desvio-padrão e da tabela completa. Uma avaliação com
divisão aleatória de janelas reportaria um número maior e menos honesto.
\end{exemplo}

\section{Pré-processamento}

As etapas de pré-processamento empregadas no livro formam um repertório pequeno.

\begin{enumerate}[leftmargin=*,itemsep=0.45em]

\item \textbf{Tratamento de ausentes}, conforme descrito acima.

\item \textbf{Normalização Min-Max} para o intervalo $[0,1]$:
\[
  \tilde{x} = \frac{x - \min(x)}{\max(x) - \min(x)}.
\]
Usada no \cref{cap:dual}, onde o autoencoder tem saída sigmoide. Sensível a
\ing{outliers} e ao intervalo de calibração.

\item \textbf{Padronização} (escore $z$):
\[
  \tilde{x} = \frac{x - \mu}{\sigma}.
\]
Usada nos \cref{cap:segmentacao,cap:mahalanobis}. Preferível quando o modelo assume
distribuição aproximadamente gaussiana --- caso, notadamente, da distância de
Mahalanobis.

\item \textbf{Filtragem por transformada \ing{wavelet}.} Decomposição com \ing{wavelet}
de Daubechies db4 em três níveis, aplicação de limiar suave aos coeficientes de detalhe
e reconstrução. Remove ruído de alta frequência preservando descontinuidades --- o que
um filtro passa-baixas convencional não faz. É particularmente relevante aqui, já que
as assinaturas de interesse \emph{são} descontinuidades.

\item \textbf{Injeção de ruído} de 0 a 0,5\,\% durante o treinamento, empregada no
\cref{cap:dual} como regularização e para tornar o autoencoder robusto à variabilidade
de instrumentação entre poços.

\end{enumerate}

\begin{destaque}[O conjunto de ajuste importa]
Aplicar padronização e, em seguida, normalização Min-Max não produz dois efeitos
distintos quando ambas são ajustadas sobre as mesmas amostras. Nesse caso, a
transformação afim da primeira etapa é cancelada pela segunda e o resultado coincide
com a normalização Min-Max direta. A composição só acrescenta informação quando os
parâmetros são estimados sobre referências diferentes --- por exemplo, média e desvio
do regime de base de cada poço, seguidos de limites globais estimados exclusivamente no
conjunto de treinamento. Por isso, uma implementação reprodutível deve registrar não
apenas a ordem, mas também os dados usados para ajustar cada transformação.
\end{destaque}

\begin{destaque}[Prática]
Todo o percurso deste capítulo --- carga dos arquivos do 3W, comparação entre os quatro
métodos de escalonamento, janelamento e partição por validação cruzada --- está
implementado e executável no recurso descrito no \cref{cap:pratica}, no caderno
\texttt{1\_data\_treatment} \citep{Lopes2026Intro3W}.
\end{destaque}

\resumocapitulo{%
O conjunto 3W é a referência pública da área, com \num{1984} instâncias e oito classes
de anomalia na versão 1.0 e \num{2228} instâncias e nove classes na versão 2.0. Suas
instâncias têm três procedências --- real, simulada e desenhada à mão --- e apenas as
reais sustentam conclusões operacionais. A rotulação por amostra em três estados
(normal, transiente, anômalo) reconhece que eventos têm início gradual, e o tratamento
dado ao estado transiente é uma decisão de projeto com efeito direto sobre a detecção
precoce. As principais armadilhas são o desbalanceamento severo, os sensores congelados
que passam por filtros de nulos, o vazamento temporal e entre poços em protocolos de
avaliação descuidados e, sobretudo, a correlação espúria entre início de anomalia e
perda de telemetria de fundo --- um artefato de coleta capaz de produzir modelos com
alto desempenho e nenhum valor.}

\begin{exercicios}
\item Uma instância do 3W tem \num{86400} amostras (um dia a \SI{1}{\hertz}). Calcule
      quantas janelas de $T = 200$ com passo $h = 100$ ela produz. Em seguida, calcule a
      fração de amostras compartilhadas entre duas janelas consecutivas e discuta o
      impacto disso sobre uma divisão aleatória treino/teste.

\item Implemente um detector de sensor congelado: dado um sinal e uma janela, retorne
      as regiões em que o desvio-padrão fica abaixo de $\epsilon$. Aplique-o a
      instâncias reais do 3W e reporte a fração de amostras afetadas por classe.
      Discuta se essa fração é uniforme entre classes e o que uma não uniformidade
      implicaria.

\item Compare experimentalmente três políticas de tratamento do estado transiente
      (descartar, tratar como anômalo, classe própria) em um classificador binário
      simples para a classe 2. Reporte acurácia balanceada e a distância temporal média
      entre o início real do evento e a primeira detecção.

\item Considere o achado da \cref{sec:dados-telemetria}. Proponha um experimento
      controlado que quantifique a contribuição da perda de telemetria para o
      desempenho de um classificador. Dica: pense em uma ablação em que a informação de
      ``dado ausente'' é deliberadamente removida.

\item Justifique matematicamente por que a normalização Min-Max é sensível a
      \ing{outliers} e a padronização é sensível à hipótese de simetria. Para cada uma,
      construa um exemplo numérico de dez pontos em que ela falha.

\item Projete um protocolo de validação cruzada para o 3W que simultaneamente:
      (a) evite vazamento entre poços; (b) mantenha a estratificação por classe;
      (c) garanta que cada classe apareça em pelo menos três dobras. Discuta se as três
      condições são sempre simultaneamente satisfazíveis dado o desbalanceamento do
      conjunto.
\end{exercicios}

\begin{leituras}
\item \citet{VARGAS2019} --- artigo original do 3W, com a descrição do processo de
      rotulação.
\item \citet{VazVargas2026_3w2} --- versão 2.0, com discussão explícita das limitações da
      versão anterior.
\item \citet{marins2020fault} e \citet{Marins2021} --- primeiros trabalhos de referência
      sobre o 3W, úteis como linha de base e como ilustração das convenções de avaliação
      que se consolidaram na área.
\item \citet{carvalho2021flow} e \citet{bronstad2021data} --- perspectivas
      complementares sobre qualidade de dados em séries temporais industriais.
\end{leituras}
